\clearpage\SourceLeaf{159}%
\thispagestyle{empty}
\pdfbookmark[0]{Besprechungen eigener Schriften}{notices-title}
\begin{center}
\vspace*{1in}{\Large BESPRECHUNGEN\par}
\vspace{.4in}{\large EIGENER\par}
\vspace{.4in}{\Large SCHRIFTEN.\par}
\end{center}

\clearpage\SourceLeaf{161}%
\renewcommand{\BookRunningTitle}{Besprechungen eigener Schriften}
\pdfbookmark[1]{Theorematis arithmetici demonstratio nova. Besprechung 1808 Mai 12}{notice-1808-05-12}
\begin{center}
\rule{.86\linewidth}{.3pt}\par
Göttingische gelehrte Anzeigen.\quad 1808 Mai 12.\par
\rule{.86\linewidth}{.3pt}
\end{center}

Eine vom Herrn Professor \textsc{Gauss} am 15. Januar dieses Jahres der Königlichen Gesellschaft der Wissenschaften vorgelegte Abhandlung,
\begin{center}\emph{Theorematis arithmetici demonstratio nova},\end{center}
deren Inhalt wir hier noch vorstellen müssen, behandelt den berühmten Fundamentalsatz der Lehre von den quadratischen Resten, der sowohl in der gesamten \emph{höheren Arithmetik} als auch in den angrenzenden Teilen der Analysis eine so wichtige Rolle spielt. Bekanntlich heißt eine ganze Zahl $a$ \emph{quadratischer Rest} der ganzen Zahl $b$, wenn es Zahlen der Form $xx-a$ gibt, die durch $b$ teilbar sind; im entgegengesetzten Fall heißt $a$ \emph{quadratischer Nichtrest} von $b$. Die Zahl $a$ kann positiv oder negativ sein, $b$ wird dagegen stets als positiv angesehen. Die höhere Arithmetik lehrt, dass alle Primzahlen $b$, für die eine gegebene Zahl $a$ quadratischer Rest ist, in bestimmten linearen Formen enthalten sind, während andere lineare Formen alle Primzahlen enthalten, von denen $a$ Nichtrest ist. So ist zum Beispiel $-1$ quadratischer Rest aller Primzahlen der Form $4n+1$ und quadratischer Nichtrest aller Primzahlen der Form $4n+3$; ferner ist $+2$ quadratischer Rest aller Primzahlen der Formen $8n+1$, $8n+7$, dagegen quadratischer Nichtrest aller Primzahlen der Formen $8n+3$, $8n+5$. Es gibt unendlich viele ähnliche spezielle Sätze; sie alle lassen sich aber aus der Verbindung der beiden angeführten
\SourceLeaf{162}%
mit dem folgenden allgemeinen Satz ableiten: Zwei verschiedene positive (ungerade) Primzahlen, $p$, $q$, stehen stets in \emph{derselben} wechselseitigen Beziehung (das heißt, die eine ist quadratischer Rest oder Nichtrest der anderen, je nachdem die andere Rest oder Nichtrest der ersten ist), wenn entweder beide die Form $4n+1$ haben oder wenigstens eine diese Form hat. Ihre wechselseitige Beziehung ist dagegen entgegengesetzt (das heißt, die eine ist Nichtrest der anderen, wenn diese Rest von jener ist, und umgekehrt), wenn \emph{beide} zugleich die Form $4n+3$ haben. Das ist der erwähnte Fundamentalsatz, den man in mehr als einer Gestalt ausdrücken kann. Die hier gewählte ist diejenige, in der er in der Abhandlung des Herrn Professor \textsc{Gauss} neu bewiesen wird.

Die schönsten Sätze der höheren Arithmetik, insbesondere auch die hier behandelten, haben die Eigenart, dass sie anhand von Zahlenbeispielen leicht entdeckt werden können, ihre Beweise dagegen äußerst tief verborgen liegen und nur durch sehr tiefgehende Untersuchungen aufzuspüren sind. Gerade das verleiht der höheren Arithmetik jenen zauberischen Reiz, der sie zum Lieblingsgebiet der führenden Mathematiker gemacht hat, ganz abgesehen von ihrem unerschöpflichen Reichtum, in dem sie alle anderen Teile der reinen Mathematik so weit übertrifft. Die beiden oben erwähnten speziellen Sätze waren schon \textsc{Fermat} bekannt, der nach eigener Aussage auch ihre Beweise besaß. Ob er sich darin täuschte, können wir nicht entscheiden, da er nie etwas darüber veröffentlicht hat. Für möglich dürfen wir es aber gewiss halten, da es mehrere Beispiele von Selbsttäuschung bei anderen großen Mathematikern gibt, insbesondere bei \textsc{Euler}, \textsc{Legendre} und auch bei \textsc{Fermat} selbst. Den ersten Beweis des ersten dieser Sätze gab \textsc{Euler}. Den anderen zu beweisen gelang diesem großen Mathematiker jedoch trotz seiner eifrigen, über viele Jahre fortgesetzten Bemühungen nicht; erst \textsc{Lagrange} war es vorbehalten, diese Lücke zu schließen. Beide Mathematiker bewiesen auch verschiedene andere spezielle Sätze; eine größere Anzahl aber, die sie anhand von Zahlenbeispielen fanden, entzog sich stets ihren Beweisversuchen. Es ist indessen ein merkwürdiges Spiel des Zufalls, dass beide Mathematiker auf diesem Weg nicht den allgemeinen Fundamentalsatz entdeckten, der sich so einfach formulieren lässt. Zuerst hat ihn, wenn auch in etwas anderer Gestalt, \textsc{Legendre} in der \emph{Histoire de l'Académie des Sciences de Paris} 1785 vorgetragen. Sowohl dort als auch später in seinem Werk \emph{Essai d'une théorie des nombres} versuchte
\SourceLeaf{163}%
dieser ausgezeichnete Analyst, den Beweis auf sehr scharfsinnige Untersuchungen zu gründen. Diese führten dennoch nicht zum gewünschten Ziel, das, wenn wir uns nicht irren, auf diesem Weg auch nicht erreicht werden konnte.

Der Verfasser der Abhandlung, der diese Besprechung gewidmet ist, betrat das Gebiet der höheren Arithmetik zu einer Zeit, als ihm alle früheren Arbeiten anderer Mathematiker auf diesem Gebiet völlig unbekannt waren. Diesem Umstand ist hauptsächlich zuzuschreiben, dass er überall einen ganz eigenen Weg eingeschlagen hat. Jenen Fundamentalsatz fand er zwar schon sehr früh anhand von Zahlenbeispielen; doch erst ein ganzes Jahr später gelang es ihm nach vielen Schwierigkeiten und vergeblichen Versuchen, den ersten vollkommen strengen Beweis zu finden, der im vierten Abschnitt seiner \emph{Disquisitiones Arithmeticae} entwickelt wird. Dieser Beweis beruht jedoch auf sehr mühsamen und weitläufigen Ausführungen. Später fand er noch drei andere Beweise, die zwar von dieser Unbequemlichkeit frei sind, dafür aber andere sehr tiefgehende und ihrem Inhalt nach völlig andersartige Untersuchungen voraussetzen. Einer dieser Beweise ist ebenfalls im angeführten Werk in Art.~262 mitgeteilt; die beiden anderen werden zu gegebener Zeit bekannt gemacht werden. Es blieb also immer der Wunsch, einen kürzeren Beweis zu finden, der von andersartigen Untersuchungen unabhängig wäre. Der Verfasser hofft daher, dass die Freunde der höheren Arithmetik mit Vergnügen einen fünften Beweis aufnehmen werden, der in der vorliegenden Abhandlung auf weniger als fünf Seiten dargestellt wird und in jeder Hinsicht nichts zu wünschen übrig zu lassen scheint. Bei der gedrängten Kürze dieses Beweises können wir hier freilich nur eine unvollkommene Vorstellung von seinem Verlauf geben. Mehr wäre hier umso überflüssiger, als der Band XVI der \emph{Commentationes}, in dem er bereits abgedruckt ist, demnächst erscheinen wird.

\Needspace{30\baselineskip}
Die Grundlage des Beweises ist der folgende neue Satz: Es sei $p$ eine (positive ungerade) Primzahl und $k$ eine beliebige ganze Zahl, die durch $p$ nicht teilbar ist. Wenn ferner unter den Resten, die bei der Division der $\frac12(p-1)$ Produkte $k$, $2k$, $3k\dotsfour\frac12(p-1)k$ durch $p$ entstehen, insgesamt $\mu$ größer als $\frac12p$ sind (also $\frac12(p-1)-\mu$ kleiner als $\frac12p$), so ist $k$ quadratischer Rest von $p$, wenn $\mu$ gerade ist, dagegen quadratischer Nichtrest, wenn $\mu$ ungerade ist. Die Zahl $\mu$, die nur von $k$
\SourceLeaf{164}%
und $p$ abhängt, bezeichnen wir mit dem Zeichen $(k,p)$. Durch eine Folge von Schlüssen, die sich nicht auszugsweise wiedergeben lassen, wird nun gezeigt, dass für zwei teilerfremde ungerade Zahlen $k$ und $p$ die Zahl $(k,p)+(p,k)+\frac14(k-1)(p-1)$ stets \emph{gerade} ist. Daraus folgt, dass immer dann, wenn $k$ und $p$ beide die Form $4n+3$ haben, eine der Zahlen $(k,p)$, $(p,k)$ gerade und die andere ungerade sein muss. In allen anderen Fällen dagegen, also wenn beide Zahlen $k$ und $p$ oder wenigstens eine von ihnen die Form $4n+1$ haben, müssen entweder .$(k,p)$, $(p,k)$ beide zugleich gerade oder beide zugleich ungerade sein. Zusammen mit dem obigen Satz ergibt sich daraus der Fundamentalsatz unmittelbar.---Auf demselben Weg, auf dem diese Ergebnisse gefunden werden, wird in der Abhandlung zugleich ein neuer Beweis der beiden oben erwähnten speziellen Sätze gegeben: Man zeigt nämlich leicht, dass $(-1,p)=\frac12(p-1)$ gilt, also eine gerade oder ungerade Zahl ist, je nachdem $p$ die Form $4n+1$ oder $4n+3$ hat; ebenso gilt $(2,p)=\frac14(p-1)$, wenn $p$ die Form $4n+1$ hat, und $(2,p)=\frac14(p+1)$, wenn $p$ die Form $4n+3$ hat. Daher ist $(2,p)$ gerade, wenn $p$ die Form $8n+1$ oder $8n+7$ hat, dagegen ungerade, wenn $p$ die Form $8n+3$ oder $8n+5$ hat.
\BookRule{.28\linewidth}
